\documentclass{article}
\usepackage{amsmath,amssymb,amsthm}
\usepackage{algorithm}
\usepackage{algpseudocode}
\usepackage{geometry}
\geometry{margin=1in}
\newtheorem{theorem}{Theorem}
\newtheorem{lemma}{Lemma}
\newtheorem{assumption}{Assumption}
\newcommand{\E}{\mathbb{E}}
\newcommand{\R}{\mathbb{R}}

\begin{document}

%%%%%%%%%%%%%%%%%%%%%%%%%%%%%
\section*{Quantized Stochastic Gradient Descent (QSGD)}
%%%%%%%%%%%%%%%%%%%%%%%%%%%%%

\section*{Mathematical Formulation}
Let $f:\R^{d}\rightarrow\R$ be a convex and $L$‑smooth objective.  
At iteration $t$ a (possibly stochastic) gradient $g_t$ is computed,
\[
g_t:=\nabla f(w_t;\xi_t),
\]
where $\xi_t$ denotes the random data sample (or mini‑batch).  
The algorithm does **not** transmit $g_t$ directly; instead it transmits a
quantized version $Q(g_t)$ satisfying  

\[
\E[\,Q(g_t)\mid w_t\,]=g_t ,\qquad 
\E\!\big[\|Q(g_t)-g_t\|^{2}\mid w_t\big]\le \sigma_{q}^{2}.
\]

With a (possibly time‑varying) stepsize $\eta_t>0$, the update rule is  

\[
\boxed{ \; w_{t+1}=w_t-\eta_t\,Q(g_t)\; } \tag{1}
\]

The algorithm stops after $T$ iterations or when a prescribed tolerance is met.

\section*{Pseudocode}
\begin{algorithm}[H]
\caption{Quantized Stochastic Gradient Descent}
\begin{algorithmic}[1]
\State \textbf{Input:} initial point $w_0\in\R^{d}$, stepsize schedule $\{\eta_t\}_{t\ge0}$, number of iterations $T$.
\For{$t=0,\dots,T-1$}
    \State Sample data $\xi_t$ and compute stochastic gradient $g_t=\nabla f(w_t;\xi_t)$.
    \State Quantize: $q_t = Q(g_t)$  \Comment{unbiased: $\E[q_t|w_t]=g_t$}
    \State Update: $w_{t+1}= w_t - \eta_t q_t$.
\EndFor
\State \textbf{Output:} $w_T$ (or average $\bar w_T=\frac1T\sum_{t=0}^{T-1}w_t$).
\end{algorithmic}
\end{algorithm}

\section*{Dry Run Example}
Consider the one‑dimensional quadratic $f(w)=(w-3)^2$, which is convex and $L=2$ smooth.
We use a deterministic gradient (no stochasticity) and a simple 2‑level quantizer
\[
Q(g)=\begin{cases}
+1 &\text{if } g\ge 0,\\
-1 &\text{if } g<0,
\end{cases}
\]
which is unbiased for this example because $g\in\{-2, -4\}$ etc. (the expectation is taken over the deterministic gradient, so $Q(g)=g$ in this toy case).  
Set $w_0=0$, stepsize $\eta=0.1$, and run three iterations.

\begin{center}
\begin{tabular}{c|c|c|c}
$t$ & $g_t=\nabla f(w_t)$ & $Q(g_t)$ & $w_{t+1}=w_t-\eta Q(g_t)$ \\ \hline
0 & $2(w_0-3) = -6$ & $-1$ & $0-0.1(-1)=0.1$ \\
1 & $2(w_1-3)=2(0.1-3)=-5.8$ & $-1$ & $0.1-0.1(-1)=0.2$ \\
2 & $2(w_2-3)=2(0.2-3)=-5.6$ & $-1$ & $0.2-0.1(-1)=0.3$
\end{tabular}
\end{center}

Objective values:
\[
f(w_0)=9,\quad f(w_1)= (0.1-3)^2=8.41,\quad
f(w_2)= (0.2-3)^2=7.84,\quad
f(w_3)= (0.3-3)^2=7.29 .
\]

The iterates move towards the optimum $w^{\star}=3$ despite the coarse quantization.

%%%%%%%%%%%%%%%%%%%%%%%%%%%%%
\section*{Assumptions}
%%%%%%%%%%%%%%%%%%%%%%%%%%%%%
\begin{assumption}[Convexity and Smoothness]\label{ass:convex}
$f:\R^{d}\to\R$ is convex and $L$‑smooth, i.e.
\[
f(y)\le f(x)+\langle\nabla f(x),y-x\rangle+\frac{L}{2}\|y-x\|^{2},
\quad\forall x,y\in\R^{d}.
\]
\end{assumption}
\begin{assumption}[Bounded Gradient]\label{ass:grad}
There exists $G>0$ such that $\| \nabla f(w;\xi) \| \le G$ almost surely for all $w$.
\end{assumption}
\begin{assumption}[Unbiased Quantization]\label{ass:unbiased}
The quantizer $Q:\R^{d}\to\R^{d}$ satisfies
\[
\E\!\big[Q(g)\mid g\big]=g,\qquad
\E\!\big[\|Q(g)-g\|^{2}\mid g\big]\le \sigma_{q}^{2}\quad\forall g\in\R^{d}.
\]
\end{assumption}
\begin{assumption}[Stepsize]\label{ass:stepsize}
The stepsizes are non‑increasing and satisfy
\[
\sum_{t=0}^{\infty}\eta_t =\infty,\qquad
\sum_{t=0}^{\infty}\eta_t^{2}<\infty .
\]
A typical choice is $\eta_t=\frac{c}{\sqrt{t+1}}$ with $c>0$.
\end{assumption}

%%%%%%%%%%%%%%%%%%%%%%%%%%%%%
\section*{Main Convergence Theorem}
%%%%%%%%%%%%%%%%%%%%%%%%%%%%%
\begin{theorem}[Convergence of QSGD]\label{thm:convergence}
Under Assumptions \ref{ass:convex}–\ref{ass:stepsize}, let $w^{\star}\in\arg\min f$.
For the averaged iterate $\bar w_T:=\frac{1}{T}\sum_{t=0}^{T-1}w_t$ we have
\[
\E\!\big[ f(\bar w_T)-f(w^{\star})\big]
\;\le\;
\frac{\|w_0-w^{\star}\|^{2}}{2\sum_{t=0}^{T-1}\eta_t}
\;+\;
\frac{L G^{2}+ \sigma_{q}^{2}}{2}\,
\frac{\sum_{t=0}^{T-1}\eta_t^{2}}{\sum_{t=0}^{T-1}\eta_t}.
\tag{2}
\]
In particular, with $\eta_t = \frac{c}{\sqrt{t+1}}$,
\[
\E\!\big[ f(\bar w_T)-f(w^{\star})\big]=
O\!\left(\frac{1}{\sqrt{T}}\right).
\]
\end{theorem}

\section*{Theoretical Properties}
\begin{itemize}
\item \textbf{Stability.} The unbiasedness of $Q$ guarantees that the expected update direction coincides with that of standard SGD; the extra variance $\sigma_{q}^{2}$ appears only in the second term of (2).
\item \textbf{Hyper‑parameter Sensitivity.} The constant $c$ in the stepsize schedule controls the trade‑off between the bias term $\|w_0-w^{\star}\|^{2}/(2\sum\eta_t)$ and the variance term. Larger $c$ speeds up early progress but inflates the variance contribution.
\item \textbf{Comparison to SGD.} Setting $\sigma_{q}^{2}=0$ recovers the classic SGD bound. Hence quantization degrades the rate only by an additive $O(\sigma_{q}^{2}/\sqrt{T})$ term.
\end{itemize}

\section*{Discussion}
The theorem shows that, despite communicating only a compressed gradient, QSGD preserves the optimal $\mathcal{O}(1/\sqrt{T})$ convergence rate for convex smooth problems. The proof hinges on the unbiasedness of the quantizer and a bounded second moment of the quantization error. Extensions to strongly convex or non‑convex settings follow by standard modifications (e.g., linear rates under strong convexity).

%%%%%%%%%%%%%%%%%%%%%%%%%%%%%
\section*{Preliminaries (Definitions and Lemmas)}
%%%%%%%%%%%%%%%%%%%%%%%%%%%%%
\begin{lemma}[Smoothness Inequality]\label{lem:smooth}
If $f$ is $L$‑smooth then for any $x,y\in\R^{d}$,
\[
f(y)\le f(x)+\langle\nabla f(x),y-x\rangle+\frac{L}{2}\|y-x\|^{2}.
\]
\end{lemma}
\begin{lemma}[Three‑Point Identity]\label{lem:three}
For any vectors $a,b,c$ and scalar $\eta>0$,
\[
\|a-\eta b-c\|^{2}= \|a-c\|^{2}-2\eta\langle b,a-c\rangle+\eta^{2}\|b\|^{2}.
\]
\end{lemma}

%%%%%%%%%%%%%%%%%%%%%%%%%%%%%
\section*{Main Proof (Step‑by‑Step Derivation)}
%%%%%%%%%%%%%%%%%%%%%%%%%%%%%
We prove Theorem \ref{thm:convergence}.  
All expectations are taken w.r.t.\ the randomness of the stochastic gradient and the quantizer, conditioned on the past iterates.

\begin{proof}
Let $w^{\star}$ be any optimal point. Using the update rule (1) and Lemma \ref{lem:three} we obtain for any $t\ge0$:
\[
\|w_{t+1}-w^{\star}\|^{2}
= \|w_{t}-w^{\star}\|^{2}
-2\eta_t\langle Q(g_t),\,w_{t}-w^{\star}\rangle
+\eta_t^{2}\|Q(g_t)\|^{2}.
\tag{3}
\]
\textbf{[Step 1]} Take conditional expectation $\E[\cdot\mid \mathcal{F}_t]$, where $\mathcal{F}_t$ denotes the sigma‑algebra generated by $\{w_0,\dots,w_t\}$.
\[
\E\!\big[\|w_{t+1}-w^{\star}\|^{2}\mid\mathcal{F}_t\big]
= \|w_{t}-w^{\star}\|^{2}
-2\eta_t\underbrace{\E\!\big[\langle Q(g_t),w_{t}-w^{\star}\rangle\mid\mathcal{F}_t\big]}_{\text{(a)}}
+\eta_t^{2}\underbrace{\E\!\big[\|Q(g_t)\|^{2}\mid\mathcal{F}_t\big]}_{\text{(b)}}.
\tag{4}
\]

\textbf{[Step 2]} (a): By unbiasedness (Assumption \ref{ass:unbiased}),
\[
\E\!\big[Q(g_t)\mid\mathcal{F}_t\big]=\E\!\big[\E[Q(g_t)\mid g_t]\mid\mathcal{F}_t\big]=\E[g_t\mid\mathcal{F}_t]=\nabla f(w_t).
\]
Hence
\[
\text{(a)} = \langle \nabla f(w_t),\,w_{t}-w^{\star}\rangle .
\tag{5}
\]

\textbf{[Step 3]} (b): Decompose $\|Q(g_t)\|^{2}= \|g_t\|^{2}+ \|Q(g_t)-g_t\|^{2}+2\langle g_t,Q(g_t)-g_t\rangle$.
Taking conditional expectation and using $\E[Q(g_t)-g_t\mid\mathcal{F}_t]=0$,
\[
\text{(b)} = \E\!\big[\|g_t\|^{2}\mid\mathcal{F}_t\big] + \E\!\big[\|Q(g_t)-g_t\|^{2}\mid\mathcal{F}_t\big]
\le G^{2}+ \sigma_{q}^{2},
\tag{6}
\]
where the inequality follows from Assumptions \ref{ass:grad} and \ref{ass:unbiased}.

\textbf{[Step 4]} Plug (5) and (6) into (4):
\[
\E\!\big[\|w_{t+1}-w^{\star}\|^{2}\mid\mathcal{F}_t\big]
\le \|w_{t}-w^{\star}\|^{2}
-2\eta_t\langle \nabla f(w_t),w_{t}-w^{\star}\rangle
+\eta_t^{2}\big(G^{2}+ \sigma_{q}^{2}\big).
\tag{7}
\]

\textbf{[Step 5]} By convexity,
\[
f(w_t)-f(w^{\star})\le \langle\nabla f(w_t),w_t-w^{\star}\rangle .
\tag{8}
\]
Rearranging (8) gives $-\,\langle\nabla f(w_t),w_t-w^{\star}\rangle\le -(f(w_t)-f(w^{\star}))$.
Insert this bound into (7):
\[
\E\!\big[\|w_{t+1}-w^{\star}\|^{2}\mid\mathcal{F}_t\big]
\le \|w_{t}-w^{\star}\|^{2}
-2\eta_t\big(f(w_t)-f(w^{\star})\big)
+\eta_t^{2}\big(G^{2}+ \sigma_{q}^{2}\big).
\tag{9}
\]

\textbf{[Step 6]} Take full expectation and sum (9) over $t=0,\dots,T-1$:
\[
\sum_{t=0}^{T-1}2\eta_t\,\E\!\big[f(w_t)-f(w^{\star})\big]
\le \|w_{0}-w^{\star}\|^{2}
-\E\!\big[\|w_{T}-w^{\star}\|^{2}\big]
+\big(G^{2}+ \sigma_{q}^{2}\big)\sum_{t=0}^{T-1}\eta_t^{2}.
\tag{10}
\]
Since the norm term is non‑negative we drop it:
\[
\sum_{t=0}^{T-1}2\eta_t\,\E\!\big[f(w_t)-f(w^{\star})\big]
\le \|w_{0}-w^{\star}\|^{2}
+\big(G^{2}+ \sigma_{q}^{2}\big)\sum_{t=0}^{T-1}\eta_t^{2}.
\tag{11}
\]

\textbf{[Step 7]} Divide both sides of (11) by $2\sum_{t=0}^{T-1}\eta_t$ and use the convexity of $f$ together with Jensen’s inequality:
\[
\E\!\big[f(\bar w_T)-f(w^{\star})\big]
\le
\frac{\|w_{0}-w^{\star}\|^{2}}{2\sum_{t=0}^{T-1}\eta_t}
+\frac{G^{2}+ \sigma_{q}^{2}}{2}\,
\frac{\sum_{t=0}^{T-1}\eta_t^{2}}{\sum_{t=0}^{T-1}\eta_t}.
\tag{12}
\]
Equation (12) is precisely the bound (2) stated in Theorem \ref{thm:convergence}.

\textbf{[Step 8]} Choose the standard diminishing stepsize $\eta_t = \frac{c}{\sqrt{t+1}}$. Then
\[
\sum_{t=0}^{T-1}\eta_t = c\sum_{t=0}^{T-1}\frac{1}{\sqrt{t+1}}
\ge 2c(\sqrt{T}-1),\qquad
\sum_{t=0}^{T-1}\eta_t^{2}=c^{2}\sum_{t=0}^{T-1}\frac{1}{t+1}
\le c^{2}(1+\log T).
\]
Plugging these estimates into (12) yields
\[
\E\!\big[f(\bar w_T)-f(w^{\star})\big]
\le
\underbrace{\frac{\|w_{0}-w^{\star}\|^{2}}{4c\sqrt{T}}}_{\displaystyle O(T^{-1/2})}
\;+\;
\underbrace{\frac{(G^{2}+ \sigma_{q}^{2})c(1+\log T)}{4c\sqrt{T}}}_{\displaystyle O(T^{-1/2})},
\]
which is $O(1/\sqrt{T})$.  ∎
\end{proof}

%%%%%%%%%%%%%%%%%%%%%%%%%%%%%
\section*{Rate Derivation (With Constants)}
%%%%%%%%%%%%%%%%%%%%%%%%%%%%%
From (12) we may write the explicit constant‑wise bound
\[
\E\!\big[f(\bar w_T)-f(w^{\star})\big]
\le
\frac{D^{2}}{2\sum_{t=0}^{T-1}\eta_t}
+\frac{(G^{2}+\sigma_{q}^{2})}{2}\,
\frac{\sum_{t=0}^{T-1}\eta_t^{2}}{\sum_{t=0}^{T-1}\eta_t},
\quad D:=\|w_{0}-w^{\star}\|.
\]
For $\eta_t = c/(t+1)^{\alpha}$ with $0.5\le\alpha<1$ the two fractions both scale as $T^{\alpha-1}$, confirming the $O(T^{-1/2})$ rate when $\alpha=1/2$.

%%%%%%%%%%%%%%%%%%%%%%%%%%%%%
\section*{Special Cases}
%%%%%%%%%%%%%%%%%%%%%%%%%%%%%
\begin{itemize}
\item \textbf{No Quantization.} $\sigma_{q}^{2}=0$ reduces (12) to the classic SGD bound.
\item \textbf{Deterministic Gradients.} If $g_t=\nabla f(w_t)$ exactly, $G$ can be replaced by $\sup_{w}\|\nabla f(w)\|$, and the same rate holds.
\item \textbf{Strong Convexity.} If $f$ is $\mu$‑strongly convex, a similar derivation gives a linear rate
\[
\E\!\big[\|w_T-w^{\star}\|^{2}\big]\le (1-\mu\eta)^{T}\|w_0-w^{\star}\|^{2}
+\frac{\eta(G^{2}+ \sigma_{q}^{2})}{\mu},
\]
showing that quantization only inflates the steady‑state error term.
\end{itemize}

\end{document}